% Eksplorator fraktali - zarys projektu do akceptacji (wersja skrocona)
% Kompilacja: pdflatex zarys_projektu_v3.tex (logo: plik AGH.pdf w tym samym folderze)

\documentclass[11pt,a4paper]{article}

\usepackage[utf8]{inputenc}
\usepackage[T1]{fontenc}
\usepackage[polish]{babel}
\usepackage{lmodern}
\usepackage[margin=2.2cm]{geometry}
\usepackage{amsmath,amssymb}
\usepackage{enumitem}
\usepackage{booktabs}
\usepackage{graphicx}
\usepackage{fancyvrb}
\usepackage[hidelinks]{hyperref}

\setlist{itemsep=3pt,topsep=4pt}
\setlength{\parskip}{5pt}
\setlength{\parindent}{0pt}
\setlength{\emergencystretch}{3em}

\newcommand{\kod}[1]{\texttt{#1}}
\newcommand{\naglowek}[1]{\par\bigskip{\large\bfseries\underline{#1}}\par\smallskip}
\fvset{frame=single,fontsize=\small,framesep=6pt}

\begin{document}

\begin{titlepage}
    \centering
    \IfFileExists{AGH.pdf}
      {\includegraphics[width=0.3\textwidth]{AGH.pdf}}
      {\fbox{\parbox[c][4cm][c]{0.3\textwidth}{\centering logo AGH}}}\\[1cm]
    \textsc{\LARGE Akademia Górniczo-Hutnicza\\
    im. Stanisława Staszica w Krakowie}\\[2cm]
    \textsc{\Large Wydział Fizyki i Informatyki Stosowanej}\\[1cm]
    \textsc{\large Interfejsy użytkownika i biblioteki graficzne}\\[1cm]
    {\huge \bfseries Eksplorator fraktali}\\[0.4cm]
    {\Large Zarys projektu do akceptacji}\\[2cm]
    {\Large \textbf{Autorzy}}\\[0.5cm]
    {\large Jakub Piela (kapitan)\\[2pt] Marcel Gładysz\\[2pt] Patryk Pilch}\\
    \vfill
    \begin{flushright}
    \textbf{Prowadzący:} dr inż. Janusz Malinowski\\
    \textbf{Temat:} własny\\
    \textbf{Rok akademicki:} 2026/2027\\
    \textbf{Data złożenia:} 02.11.2026\\
    \end{flushright}
\end{titlepage}

% ---------------------------------------------------------------
\naglowek{Opis projektu}

Celem projektu jest napisanie programu do generowania i interaktywnego
oglądania fraktali. Program obsługuje trzy rodziny fraktali:
\begin{itemize}
  \item \textbf{zbiory zespolone} -- zbiór Mandelbrota i zbiory Julii,
  \item \textbf{fraktale IFS} -- kształty definiowane zestawem przekształceń
        afinicznych, np.\ trójkąt Sierpińskiego i paproć Barnsleya,
  \item \textbf{L-systemy} -- fraktale opisane regułami przepisywania napisów
        i rysowane grafiką żółwia, np.\ krzywa Kocha i rośliny.
\end{itemize}

Użytkownik przybliża i przesuwa widok myszą, zmienia parametry i palety
kolorów w panelu sterowania, a bieżący obraz zapisuje do pliku. Program tworzy
także animację płynnego przejścia (,,morfingu'') jednego fraktala IFS w drugi
i zapisuje ją na dysku jako ciąg bitmap. Ta część nawiązuje do tematu nr~25
(,,Generowanie i animacja fraktali'').

% ---------------------------------------------------------------
\naglowek{Narzędzia}

\begin{itemize}
  \item \textbf{Język i środowisko:} C++17, Visual Studio 2026, Windows 11.
  \item \textbf{SFML 3.1} -- okno, obsługa myszy i klawiatury, wyświetlanie
        obrazu oraz zapis plików PNG i BMP.
  \item \textbf{Dear ImGui} z nakładką \textbf{ImGui-SFML} -- panel sterowania,
        menu, pola i suwaki.
  \item \textbf{ImGuiFileDialog} -- okna wyboru plików przy wczytywaniu
        i zapisie.
\end{itemize}

% ---------------------------------------------------------------
\naglowek{Wymagania podstawowe}

Program powinien realizować wymienione niżej zadania.

\begin{itemize}
  \item \textbf{Interfejs graficzny:} okno z obszarem fraktala, panel
        sterowania z polami i suwakami do zmiany parametrów, menu oraz okna
        wyboru plików. Po zmianie rozmiaru okna obraz jest przeliczany, a nie
        rozciągany.
  \item \textbf{Zbiór Mandelbrota:} przybliżanie kółkiem myszy w punkt pod
        kursorem, przesuwanie przeciąganiem, powrót do widoku startowego
        i regulowana maksymalna liczba iteracji.
  \item \textbf{Zbiory Julii:} parametr $c$ wybierany kliknięciem na zbiorze
        Mandelbrota albo wpisywany w panelu.
  \item \textbf{Kolorowanie zbiorów zespolonych:} płynne przejścia kolorów
        (bez widocznych pasów) i co najmniej trzy palety do wyboru.
  \item \textbf{Fraktale IFS:} rysowanie metodą gry w chaos, wbudowane trójkąt
        Sierpińskiego i paproć Barnsleya, edycja współczynników w panelu oraz
        wczytywanie i zapis definicji w pliku tekstowym. Kolor punktu zależy od
        przekształcenia, którym go wyznaczono. Program sam dobiera skalę tak,
        aby fraktal mieścił się w oknie i nie był zbyt mały.
  \item \textbf{Animacja przejścia:} płynne przejście między dwoma fraktalami
        IFS w zadanej przez użytkownika liczbie klatek. Program pokazuje klatki
        w trakcie generowania i zapisuje je jako ponumerowane pliki PNG lub BMP.
  \item \textbf{L-systemy:} aksjomat, reguły, kąt i liczba iteracji ustawiane
        w panelu lub wczytywane z pliku, rysowanie grafiką żółwia
        z rozgałęzieniami, wbudowane krzywa Kocha, smok Heighwaya i roślina.
        Rysunek jest dopasowywany do okna.
  \item \textbf{Zapis obrazu:} bieżący widok do pliku PNG lub BMP.
\end{itemize}

Definicja fraktala IFS to plik tekstowy, w którym każdy wiersz opisuje jedno
przekształcenie afiniczne
\[
  x' = ax + by + e, \qquad y' = cx + dy + f
\]
współczynnikami \kod{a b c d e f} oraz prawdopodobieństwem \kod{p}, z jakim
przekształcenie jest losowane w grze w chaos. Przykład -- paproć Barnsleya:

\begin{Verbatim}
 0.00  0.00  0.00  0.16  0.00  0.00  0.01
 0.85  0.04 -0.04  0.85  0.00  1.60  0.85
 0.20 -0.26  0.23  0.22  0.00  1.60  0.07
-0.15  0.28  0.26  0.24  0.00  0.44  0.07
\end{Verbatim}

Klatki animacji przejścia od fraktala $A$ do fraktala $B$ powstają przez płynną
zmianę każdego współczynnika:
\[
  w^{(t)} = (1 - t)\, w^{A} + t\, w^{B}, \qquad t \in [0, 1] .
\]
Gdy fraktale mają różną liczbę przekształceń, do mniejszego zestawu dokładamy
kopie istniejących przekształceń.

% ---------------------------------------------------------------
\naglowek{Wymagania rozszerzone}

W wersji rozszerzonej program dodatkowo powinien:

\begin{itemize}
  \item liczyć zbiory zespolone wielowątkowo, aby przybliżanie i przesuwanie
        były płynne, a animację generować w tle z paskiem postępu i możliwością
        przerwania,
  \item tworzyć animację przez więcej niż dwa fraktale według pliku scenariusza
        (kolejne fraktale i liczba klatek każdego przejścia),
  \item rysować fraktale IFS trójwymiarowe z pozycją obserwatora ustawianą
        przez użytkownika; punkty dalsze są ciemniejsze i nie przesłaniają
        bliższych,
  \item tworzyć animację przybliżania zbioru Mandelbrota lub Julii,
  \item udostępniać edytor palet kolorów z zapisem i wczytywaniem palet,
  \item obsługiwać kolejne fraktale zespolone: Burning Ship, Multibrot ($z^n$)
        i fraktal Newtona,
  \item zapisywać obraz w rozdzielczości większej niż okno, z wygładzaniem
        krawędzi.
\end{itemize}

\end{document}